\documentclass[12pt]{letter}
\usepackage{geometry}\geometry{margin=1in}
\usepackage{enumitem}\usepackage{hyperref}

\begin{document}
\begin{letter}{Editor-in-Chief\\Expert Systems with Applications\\Elsevier}

\opening{Dear Editor,}

We submit for your consideration our manuscript titled
\textbf{``A Data-Adaptive Hybrid Quantum-Classical Transformer for Medical
Tabular Classification''} for publication in \textit{Expert Systems with
Applications}.

\textbf{Summary of contribution.}
We integrate a variational quantum circuit (VQC) into the feed-forward sublayer
of a TabTransformer encoder, and select the VQC complexity adaptively per
cross-validation fold from the training-set size. Evaluated on four medical
datasets (UCI CKD, PIMA Diabetes, Cleveland Heart Disease, Framingham Heart
Study) via 10-fold stratified cross-validation against XGBoost, LightGBM,
TabTransformer, and an MLP, the model attains the best Matthews correlation,
tied-best accuracy, and the lowest cross-fold variance on the Cleveland Heart
Disease dataset, and ranks second of five by F1 on the deliberately hard,
class-imbalanced Framingham benchmark. We report results honestly: on the
near-saturated CKD and on PIMA the model is competitive but does not lead, and
we state this plainly rather than over-claiming.

\textbf{Novel contributions.}
\begin{itemize}
  \item A data-adaptive circuit selector that scales VQC depth/width to the
        available data, with an ablation showing that on small datasets a
        shallower circuit outperforms the deeper one -- a data-regime sensitivity
        rarely reported for hybrid quantum models.
  \item Quantitative circuit characterisation -- Meyer--Wallach expressibility
        (0.995--0.9996) and entanglement capability (0.84--0.95), plus kernel
        target alignment (reported even though the quantum kernel does not beat
        an RBF kernel here, in the interest of honest benchmarking).
  \item Quantum gradient feature attribution complementing classical SHAP.
  \item Full SHA-256 cryptographic provenance for datasets and model artifacts,
        with a 45-test automated verification suite and a one-command sanity check.
\end{itemize}

\textbf{Reproducibility.}
All code, preprocessing pipelines, and results are public at
\url{https://github.com/Nafis878/-HQCT}. Running \texttt{python scripts/sanity\_check.py} verifies all reported
results against SHA-256-signed artifacts in under two minutes.

\textbf{Scope fit.}
The work sits at the intersection of quantum machine learning and clinical
decision support -- an emerging area aligned with ESWA's focus on intelligent
systems for real-world applications.

We confirm this manuscript is original, not under consideration elsewhere, and
that all authors have approved the submission.

\closing{Sincerely,}

Nafis\\
[Institution]\\
[Email]\\
[Date]

\end{letter}
\end{document}
